\subsection*{Connectome data}
\textbf{FlyWire v783.} We used the exact network distributed with the Shiu et al.\ model
(\texttt{Connectivity\_783.parquet}, \texttt{Completeness\_783.csv}; repository commit
\texttt{91bdd1e}). It has 138{,}639 neurons and 15{,}091{,}983 directed connections carrying
54{,}492{,}922 synapses, with no synapse-count threshold. Each presynaptic neuron's sign is as
given in that file. Cell types, sides and super-classes come from Schlegel et al.\
\citep{schlegel2024} (Supplemental File~1).
\textbf{BANC v888.} We used the neuron table (\texttt{banc\_888\_meta}), the connection
table built from synapse table v2 (\texttt{banc\_888\_edgelist\_simple\_v2}) and the
neuron-level transmitter predictions. We kept proofread neurons and dropped glia, giving
150{,}802 neurons, of which 126{,}389 are in the central brain or optic lobes. A neuron's
cell type is that of its FlyWire match (\texttt{fafb\_cell\_type}) when this is a FlyWire
type name, otherwise its BANC type; 91.7\% of brain neurons received a FlyWire type.
\textbf{Male CNS v1.0.} We used the connection table restricted to traced bodies with synapse
confidence $\ge0.5$
(\texttt{connectome-weights-male-cns-v1.0-minconf-0.5-traced-only}), the body annotations,
and the neuron-level transmitter predictions (\texttt{predicted\_nt}). Cell types were
harmonised to FlyWire names through the \texttt{flywireType} field, falling back to the male
CNS type. This gave 165{,}122 neurons (144{,}717 in the brain subgraph); 94.9\% of the
brain neurons have a FlyWire type.
\textbf{Signs.} As in Shiu et al., GABA and glutamate are inhibitory and all other
transmitters are excitatory. FlyWire used the signs in Shiu's file. BANC and the male CNS used
their own neuron-level predictions (``own''). Two alternative sign schemes were tested in the
robustness sweep: the same rule applied to each dataset's neuron-level prediction (``nt'';
for FlyWire, the Schlegel \texttt{top\_nt}), and the majority sign of each cell type
(``type'').
\textbf{Brain subgraph.} The primary analysis used the brain portion of each connectome,
because FlyWire has no nerve cord. For the male CNS we removed bodies whose super-class is a
nerve-cord class (\texttt{vnc\_*}, enteric, efferent-ascending) and kept descending and
ascending neurons. For BANC we kept neurons whose region is central brain or optic lobe. Whole-CNS networks were a secondary analysis.

\subsection*{Model}
We used the Shiu et al.\ model unchanged \citep{shiu2024}. Each neuron $i$ has membrane
potential $v_i$ and synaptic conductance-like variable $g_i$ (in mV):
\begin{equation}
\tau_m \frac{dv_i}{dt} = (v_0 - v_i) + g_i, \qquad \tau_g \frac{dg_i}{dt} = -g_i ,
\end{equation}
with $v_0=v_{\mathrm{reset}}=-52$\,mV, threshold $-45$\,mV, $\tau_m=20$\,ms, $\tau_g=5$\,ms,
refractory period 2.2\,ms and synaptic delay 1.8\,ms \citep{kakaria2017,shiu2024}. A spike of
neuron $j$ increments $g_i$ by $w_{ji}=\Wsyn\, s_j\, n_{ji}$ after the delay, where $n_{ji}$
is the synapse count, $s_j=\pm1$ the sign and $\Wsyn=0.275$\,mV. On a spike, $v$ and $g$ are
reset to $-52$\,mV and 0, and both are frozen while refractory. Stimulated neurons receive
Poisson input at rate $r$; each input event adds $250\,\Wsyn$ to $v$, and these neurons have
no refractory period. Silent networks remain silent: there is no background drive. Each
experiment consisted of 30 trials of 1\,s, as in Shiu et al.

\textbf{Reference implementation.} Shiu's Brian2 \citep{stimberg2019} code was run unchanged
except that the network was built once and restored before each trial (2.9.0, Cython
backend, $dt=0.1$\,ms, exact linear integration). We measured two Brian2 behaviours that
any reimplementation must reproduce. First, synaptic input arriving at a refractory neuron is
discarded rather than delayed (also reported in \citep{santos2026}). Second, within one time
step the order is state update, threshold, synaptic delivery, reset.

\textbf{Batched implementation.} We reimplemented the model in PyTorch \citep{paszke2019}.
Many trials and experiments run as rows of a state matrix. Integration is exact over each
0.1\,ms step. Synaptic delivery is event-driven: at each step the outgoing synapses of the
neurons that spiked 18 steps earlier are gathered from a compressed sparse-row matrix and
scattered into $g$, with input to refractory targets discarded. All arithmetic is in float32,
so it runs on consumer GPUs with weak double precision.

\textbf{Validation.} Before any cross-connectome run, we compared the PyTorch model (float32;
Apple-silicon GPU and CPU) with Brian2 (float64) on FlyWire v783 in three experiments: sugar
GRNs at 100\,Hz, sugar plus bitter GRNs at 100\,Hz each, and JO-CE neurons at 100\,Hz. Each
used 30 trials and independent seeds. The pre-specified pass criteria were: the per-neuron
mean-rate correlation with Brian2 is no more than 0.01 below the Brian2-vs-Brian2 seed
ceiling; the readout neuron's 95\% confidence intervals overlap; and total spike counts per
trial do not differ (Mann--Whitney $p>0.01$).

\subsection*{Stimuli and readouts}
Stimuli were defined by cell type, identically in all datasets: sugar GRNs (LB3b, LB3c),
bitter GRNs (LB1a--d), water GRNs (LB3a), Johnston's-organ neurons JO-C/E (all JO-C* and
JO-E* types) and JO-F (JO-F*). All neurons of the type(s) on one side were stimulated at
150\,Hz, Shiu's default rate. Each protocol was run on each side. Readouts were MN9 (FlyWire
type CB0701; male CNS type MN9) and aBN1 (SAD093), averaged over both sides. For the
reproduction on FlyWire we also used Shiu's exact neuron lists, whose v783 composition is
LB3c (9), LB3d (5), LB4b (4) and LB3b (2) for ``sugar''.

\subsection*{Weight normalisation}
Connectomes differ in synapse detection density, while $\Wsyn$ was set on FlyWire. We
pre-specified three normalisations, none tuned on any readout. N0: raw counts, as in Shiu et
al. N1: one global factor per dataset, $s_d = \tilde{m}_{\mathrm{FlyWire}}/\tilde{m}_d$,
where $\tilde{m}$ is the median total input synapse count over neurons of the cell types
present in all three datasets, excluding readout types. N2: input fraction, $n_{ji}$ divided
by neuron $i$'s total input and multiplied by $\tilde{m}_{\mathrm{FlyWire}}$. N1 was
designated primary.

\subsection*{Agreement metrics and ladder}
For each experiment we computed the mean rate of every (cell type, side relative to the
stimulus) group. Silent neurons count as zero, and the stimulated types are excluded. Two
profiles were compared over the groups present in both networks with Spearman's $\rho$ on
$\log(1+\mathrm{rate})$, restricted to groups active ($\ge2$\,Hz) in either profile. We also
report the Jaccard index of the active sets and the overlap of the top 20 and top 50 groups.
These choices follow \citep{santos2026}, so the numbers are directly comparable. 95\%
confidence intervals come from 1{,}000 bootstrap resamples over groups. The ladder rungs were:
\emph{seed} (the same network, two independent sets of 30 trials); \emph{hemisphere} (left vs
mirrored right stimulation in one network); \emph{between} (FlyWire--BANC, FlyWire--male CNS,
BANC--male CNS, same side); \emph{degraded} (FlyWire vs degraded FlyWire); and \emph{shuffle}
(each network vs its degree-preserving shuffles).

\subsection*{Controls}
\textbf{Degraded FlyWire.} For each cell type shared with a reference connectome (BANC or the
male CNS), let $p_t=\min(1,\tilde{m}^{\mathrm{ref}}_t/\tilde{m}^{\mathrm{FW}}_t)$, where
$\tilde{m}_t$ is the median total input synapse count of type $t$. Each connection onto a
neuron of type $t$ was thinned binomially, $n'\sim\mathrm{Bin}(n,p_t)$. Types absent from the
reference used the median ratio. This keeps FlyWire's partners, signs and individuality while
imposing the reference's per-type completeness deficit.
\textbf{Degree-preserving shuffle.} The postsynaptic partners of all connections were permuted
\citep{maslov2002}. This preserves every neuron's in- and out-degree and out-strength.
Duplicate pairs were merged and self-connections removed.
\textbf{Edge noise.} Synapse counts were resampled as $n'\sim\mathrm{Poisson}(n)$.

\subsection*{Replication criteria and statistics}
An activation prediction (sugar$\to$MN9, water$\to$MN9, JO-C/E$\to$aBN1) replicated in a
connectome if the mean readout rate was $\ge5$\,Hz and $\ge20\%$ of FlyWire's under the same
normalisation, and greater than zero by a one-sided Wilcoxon signed-rank test. Bitter
suppression replicated if MN9 during sugar plus bitter stimulation was lower than during sugar
alone (one-sided Mann--Whitney) and the suppression was at least 50\%. JO specificity
replicated if aBN1 responded more to JO-C/E than to JO-F (one-sided Mann--Whitney).
$p$-values across the non-FlyWire tests were Holm-corrected \citep{holm1979}.
After the first tables we added one condition to the last two criteria: each also requires
the underlying activation ($\ge5$\,Hz at the readout), because suppression of an absent
response is not meaningful. We also added a descriptive label, ``brain-wide activity'', for
experiments in which more than 2\% of neurons were active. Both were defined after seeing
data.
Simulation trials are the unit of replication. Confidence intervals therefore describe
simulation noise and cell-type sampling, not variation across flies: each connectome is one
animal.

\subsection*{Robustness sweep and stimulation screen}
\todo{grid as run; see src/sweep.py and notes/decisions.md D6.}
